%% kapa: exam=1 round=25 year=2014 period=2 form=A subject=appraisal_law subject_ko=부동산관계법규 seq=28 number=28
%% source: Q-Net 감정평가사 기출문제 (cst00309) artlSeq=5101207 file=extracted/1008/5101207/2014년 제25회 감정평가사 1차시험 기출문제/2교시(A형).hwp
%% conversion: kordoc 4.24.51 -> tools/gen1.py (auto); review=proofread-sonnet; flags=none
\begin{question}{28}
국유재산법상 행정재산에 관한 설명으로 옳지 않은 것은?
\begin{choices}
  \item 행정재산을 사용허가한 때에 징수하는 사용료는 선납하여야 한다.
  \item 중앙관서의 장은 그 소속 공무원에게 행정재산 관리에 관한 사무를 위임하거나 분장하게 한 경우에는 그 뜻을 국토교통부장관에게 통지하여야 한다.
  \item 중앙관서의 장은 공용ㆍ공공용ㆍ기업용 재산에 대하여 그 용도나 목적에 장애가 되지 아니하는 범위에서 사용허가를 할 수 있다.
  \item 행정재산에 대하여 주거용으로 사용허가를 하는 경우에는 수의(隨意)의 방법으로 사용허가를 받을 자를 결정할 수 있다.
  \item 중앙관서의 장은 행정재산으로 사용하기로 결정한 날부터 5년이 지난 날까지 행정재산으로 사용되지 아니한 경우에는 지체 없이 행정재산의 용도를 폐지한다.
\end{choices}
\end{question}
\answer{2}
